% Part: applied-modal-logics
% Chapter: epistemic-logic
% Section: language-epistemic-logic

\documentclass[../../../include/open-logic-section]{subfiles}

\begin{document}

\olfileid{aml}{el}{lan}

\olsection{د معرفتي منطق ژبه}

\begin{defn}
فرض کړئ $G$ د عاملانو د نښو يو سټ وي۔ د څو-عامله معرفتي منطق
بنسټيزه ژبه دا توکي لري:
\begin{enumerate}
  \tagitem{prvFalse}{د !!{falsity} قضيوي ثابت~$\lfalse$۔}{}
  \tagitem{prvTrue}{د !!{truth} قضيوي ثابت~$\ltrue$۔}{}
  \item د !!{propositional variable}s يو !!{denumerable}s سټ: $\Obj
    p_0$، $\Obj p_1$، $\Obj p_2$، \dots
  \item قضيوي نښلوونکي: \startycommalist
  \iftag{prvNot}{\ycomma $\lnot$ (نفي)}{}%
  \iftag{prvAnd}{\ycomma $\land$ (عطف)}{}%
  \iftag{prvOr}{\ycomma $\lor$ (فصل)}{}%
  \iftag{prvIf}{\ycomma $\lif$ (!!{conditional})}{}%
  \iftag{prvIff}{\ycomma $\liff$ (!!{biconditional})}{}.
  \item د پوهې عامل $\Knows_a$، چې $a \in G$ وي۔
\end{enumerate}
\end{defn}

که په خپل نظام کښې يوازې د يوه عامل پوهه څېړو، د سټ~$G$ او د
جلا جلا عاملانو يادونه پرېښودلے شو۔ په دې حالت کښې يوازې بنسټيز
عامل~$\Knows$ لرو۔

\begin{defn}
د معرفتي ژبې \emph{!!^{formula}s} په استقرايي ډول داسې
  تعريفېږي:
\begin{enumerate}
\tagitem{prvFalse}{$\lfalse$ يو اټومي !!{formula} دے۔}{}

\tagitem{prvTrue}{$\ltrue$ يو اټومي !!{formula} دے۔}{}

\item هر قضيوي متحول~$\Obj p_i$ يو اټومي !!{formula} دے۔

\tagitem{prvNot}{که $!A$ يو !!a{formula} وي، نو $\lnot !A$ هم
  يو !!a{formula} دے۔}{}

\tagitem{prvAnd}{که $!A$ او $!B$ !!{formula}s وي، نو $(!A \land
  !B)$ يو !!a{formula} دے۔}{}

\tagitem{prvOr}{که $!A$ او $!B$ !!{formula}s وي، نو $(!A \lor !B)$
  يو !!a{formula} دے۔}{}

\tagitem{prvIf}{که $!A$ او $!B$ !!{formula}s وي، نو $(!A \lif !B)$
  يو !!a{formula} دے۔}{}

\tagitem{prvIff}{که $!A$ او $!B$ !!{formula}s وي، نو $(!A \liff !B)$
  يو !!a{formula} دے۔}{}

\item که $!A$ يو !!a{formula} وي او $a \in G$ وي، نو $\Knows_a !A$
  هم يو !!a{formula} دے۔

\tagitem{limitClause}{بل هېڅ شے !!a{formula} نۀ دے۔}{}
\end{enumerate}
\end{defn}

که !!a{formula}~$!A$ کښې $\Knows_a$ نۀ وي، هغه
\emph{له موداليت پرته} بولو۔

\begin{defn}
  د $\Knows$ عامل د يوه کس د پوهې د ښودلو لپاره دے؛ $\EKnows$،
  چې ډېر ځله «هر څوک پوهېږي» لوستل کېږي، د ډلې د هر غړي پوهه ښيي۔
  که $G' \subseteq G$ وي، نو $\EKnows_{G'} !A$ د
  \[\bigwedge_{b \in G'} \Knows_b !A.\] لنډه نښه تعريفوو۔
\end{defn}

د عاملانو د ډلې~$G$ په منځ کښې د پوهې يوه لا پياوړې معنا هم
تعريفولے شو، يعنې \emph{مشترکه پوهه}۔ که د معلوماتو يوه برخه د
ډلې مشترکه پوهه وي، نو د ډلې د عاملانو په هر ممکن ترکيب کښې هر يو
پوهېږي چې نور پوهېږي، او پوهېږي چې نور پوهېږي چې نور پوهېږي،
\dots بې‌پايه۔ دا د ډلې د هر غړي تر پوهې ډېره پياوړې ده۔ داسې
اړيکيز مدلونه جوړول اسانه دي چې پکښې يو !!a{formula} د ډلې د هر
غړي پوهه وي، خو مشترکه پوهه نۀ وي۔ $\CKnows_G !A$ به د دې
ښودلو لپاره کاروو چې «په~$G$ کښې دا مشترکه پوهه ده چې~$!A$»۔

\end{document}
